\section*{Chapter \ref{ch:m1:shares-corporate-actions-indices} --- Shares, Corporate Actions and Indices}

\begin{solution}{exo:m1:shares-corporate-actions-indices:1}
\$52 billion; float-adjusted \$36.4 billion.
\end{solution}

\begin{solution}{exo:m1:shares-corporate-actions-indices:2}
Each share becomes 1.5 shares: opening price $240/1.5 = \$160$, factor
$2/3$, position 750 shares (still worth \$120\,000).
\end{solution}

\begin{solution}{exo:m1:shares-corporate-actions-indices:3}
Ex price \$78.80; $f = 78.80/80.00 = 0.985$; $76.00 \times 0.985 = \$74.86$.
\end{solution}

\begin{solution}{exo:m1:shares-corporate-actions-indices:4}
$\mathrm{TERP} = (7\times25 + 2\times16)/9 = 23.00$; one right is worth
$23 - 16 = 7.00$; $f = 0.92$. A holder of 700 shares receives rights to 200
new shares, sells them for \euro 1\,400, and holds $700 \times 23 = 16\,100$:
\euro 17\,500, as before.
\end{solution}

\begin{solution}{exo:m1:shares-corporate-actions-indices:5}
Level $500/0.5 = 1\,000$. After the split the third price is 110 and the sum
280, so $\Delta' = 280/1\,000 = 0.28$. Weights go from 10\%, 24\%, 66\% to
17.9\%, 42.9\%, 39.3\%: the index is now a different portfolio, by a decision
of the company's board and not of the index provider.
\end{solution}

\begin{solution}{exo:m1:shares-corporate-actions-indices:6}
Standard factor: $48.45/47.50 - 1 = +2.00\%$. Total return: $(48.45 +
2.50)/50 - 1 = +1.90\%$. Difference 0.10\%, equal to $\varepsilon D/P = 0.02
\times 0.05$.
\end{solution}

\begin{solution}{exo:m1:shares-corporate-actions-indices:7}
(a) Market value \$5 billion, divisor 5\,000\,000. (b) Market value falls to
$2.0 + 2.4 = \$4.4$ billion at unchanged prices: divisor 4\,400\,000, level
still 1\,000. (c) $(100\times22 + 40\times57)/4.4 = 1\,018.2$.
\end{solution}

\begin{solution}{exo:m1:shares-corporate-actions-indices:8}
First, the adjusted price of a stock twenty years ago is its raw price times
every later factor: a stock that has since split many times shows an old
adjusted price below \$5 although it traded at \$80 that day. The rule
therefore selects, in the past, precisely the stocks that later rose enough
to split: the future is in the filter. Second, a downloaded list of today's
3\,000 stocks contains only survivors; the low-priced stocks that went to
zero are absent. Both errors flatter the rule; the test needs raw prices as
of each date and a universe as of each date.
\end{solution}

\begin{solution}{pb:m1:shares-corporate-actions-indices:1}
\textbf{1.} P 20.0, Q 24.0, R 6.0, S 27.0: \$77.0 billion.
\textbf{2.} $77\times10^9/2\,500 = 30\,800\,000$.
\textbf{3.} 26.0\%, 31.2\%, 7.8\%, 35.1\%.
\textbf{4.} $0.01 \times 24\times10^9/30.8\times10^6 = 7.8$ points.
\textbf{5.} Price 50.00, shares 600 million; the capitalisation and hence
the divisor are unchanged.
\textbf{6.} Market value falls by $500 \times 1.00 = \$0.5$ billion: the
index is $76.5\times10^9/30.8\times10^6 = 2\,483.77$, down 16.23 points. No
adjustment: in a price index an ordinary \omterm{def:m1:shares-corporate-actions-indices:dividend}{dividend} is a price move.
\textbf{7.} It adds the 16.23 points back: 2\,500.00.
\textbf{8.} The \omterm{def:m1:shares-corporate-actions-indices:tri}{total return index} (net of withholding tax if the fund
suffers it); against the price index the fund would appear to outperform by
every \omterm{def:m1:shares-corporate-actions-indices:dividend}{dividend}.
\textbf{9.} $\mathrm{TERP} = (4\times12 + 8)/5 = 11.20$; 1\,250 million
shares.
\textbf{10.} $11.20 \times 1\,250 \times 0.5 = \$7.0$ billion; market value
$76.5 - 6.0 + 7.0 = \$77.5$ billion.
\textbf{11.} $\Delta' = 30.8 \times 77.5/76.5 = 31\,202\,614$; the level
stays 2\,483.77.
\textbf{12.} R's weight rises because the company is larger: the fund
subscribes its rights (or sells some rights and buys shares to the same
effect), investing new money in R financed by selling a little of everything
else.
\textbf{13.} $77.5 - 27.0 + 13.5 = \$64.0$ billion; $\Delta_{\text{new}} = 31\,202\,614
\times 64/77.5 = 25\,767\,320$.
\textbf{14.} P 30.5\%, Q 37.5\%, R 10.9\%, T 21.1\%.
\textbf{15.} S is $27/77.5 = 34.8\%$ of the fund: it receives \$697 million
in cash from the acquirer and must invest it across the new index, of which
\$422 million in T; the remaining \$275 million is spread over P, Q and R.
\textbf{16.} $422/150 = 2.8$ days of volume; cost $0.7\times0.02\times
\sqrt{2.8} = 2.3\%$, about \$9.9 million --- for one fund.
\textbf{17.} Traders who bought T when the change was announced, and
arbitrageurs who buy during the week in order to sell at the effective
close: the fund's cost is their revenue (One Quant Book 8).
\textbf{18.} Trackers need time to arrange the trade, and an unannounced
change would let the provider's own leak be the only information: the delay
trades a known cost (front-running) against fairness and orderly execution.
Some providers use longer notice and staggered implementation to spread the
impact.
\textbf{19.} \textbf{25\,767\,320.}
\textbf{20.} Change the level of the index.
\end{solution}

\begin{solution}{iq:m1:shares-corporate-actions-indices:1}
A 2-for-1 split (or a large special dividend or \omterm{def:m1:shares-corporate-actions-indices:spinoff}{spin-off}). The pipeline must
ingest corporate actions from a reference source, keep raw prices and
actions separately, serve back-adjusted series for return computation, and
adjust open positions, open orders and reference prices at the same moment.
A large overnight move with no action on file should raise an alert before
it reaches a strategy.
\iqlookfor{positions and orders, not only price history.}
\end{solution}

\begin{solution}{iq:m1:shares-corporate-actions-indices:2}
By the \omterm{def:m1:shares-corporate-actions-indices:dividend}{dividend}, absent taxes and news, by no-arbitrage between the cum close
and the ex open. By less when \omterm{def:m1:shares-corporate-actions-indices:dividend}{dividends} are taxed more than capital gains
for the marginal investor, and measured drops are also blurred by the
overnight market move and by the tick size.
\iqlookfor{the arbitrage argument, then the tax clientele effect.}
\end{solution}

\begin{solution}{iq:m1:shares-corporate-actions-indices:3}
Raw for anything the exchange sees: limit prices, tick rules, fill
simulation, fees. Adjusted for returns and indicators across ex-dates. Raw
used for returns: a split appears as a $-90\%$ day and the risk model
explodes. Adjusted used for orders or filters: a backtest ``buys'' at prices
that never traded, or selects stocks using information from later splits.
\iqlookfor{the look-ahead in adjusted prices.}
\end{solution}

\begin{solution}{iq:m1:shares-corporate-actions-indices:4}
A capitalisation weight is price times shares, which a split leaves
unchanged. A price weight is the price alone, which the split divides; the
divisor restores the level but cannot restore the weights.
\iqlookfor{one sentence each; bonus for ``the divisor fixes the level, not
the composition''.}
\end{solution}

\begin{solution}{iq:m1:shares-corporate-actions-indices:5}
Store: raw prices and volumes as traded, immutable; corporate actions with
ex-date, terms and announcement timestamp; identifier history. Compute on
read: \omterm{def:m1:shares-corporate-actions-indices:factor}{adjustment factors} and adjusted series for a given as-of date, so that
point-in-time queries are possible and a late correction to one action fixes
every series at once. Never store adjusted prices as primary data: they go
stale with each new action and hide the convention used.
\iqlookfor{as-of queries and immutability.}
\end{solution}

\begin{solution}{iq:m1:shares-corporate-actions-indices:6}
No. The mechanical effect of a 1-for-3 issue at a 40\% discount is a TERP
10\% below the cum price, and it appears on the \emph{ex-date}, with the
holder compensated by the rights. A fall on the \emph{announcement} is
information: the company needs capital, perhaps urgently, and management
chose to issue equity at this price; holders who cannot subscribe must sell
rights into a market that knows it. A 12\% fall on announcement is a loss of
value; the later 10\% ex-date drop is not.
\iqlookfor{separating the announcement effect from the ex-rights
adjustment.}
\end{solution}
